%%%%%%%%%%%%%%%%%%%%%%%%%%%%%%%%%%%%%%%%%
% "ModernCV" CV and Cover Letter
% LaTeX Template
% Version 1.1 (9/12/12)
%
% This template has been downloaded from:
% http://www.LaTeXTemplates.com
%
% Original author:
% Xavier Danaux (xdanaux@gmail.com)
%
% License:
% CC BY-NC-SA 3.0 (http://creativecommons.org/licenses/by-nc-sa/3.0/)
%
% Important note:
% This template requires the moderncv.cls and .sty files to be in the same 
% directory as this .tex file. These files provide the resume style and themes 
% used for structuring the document.
%

% Zijian_(Richard)_Meng_Resume_v2025-07.txt
%%%%%%%%%%%%%%%%%%%%%%%%%%%%%%%%%%%%%%%%%

%----------------------------------------------------------------------------------------
%	PACKAGES AND OTHER DOCUMENT CONFIGURATIONS
%----------------------------------------------------------------------------------------



\documentclass[12pt,letterpaper,sans]{moderncv} % Font sizes: 10, 11, or 12; paper sizes: a4paper, letterpaper, a5paper, legalpaper, executivepaper or landscape; font families: sans or roman

\moderncvstyle{classic} % CV theme - options include: 'casual' (default), 'classic', 'oldstyle' and 'banking'
\moderncvcolor{blue} % CV color - options include: 'blue' (default), 'orange', 'green', 'red', 'purple', 'grey' and 'black'


\usepackage{fontawesome5} % Needed for the GitHub icon

\makeatletter
% 1. Define the \github command taking 2 arguments (URL and Display Name)
\newcommand*{\github}[2]{\def\@github{#1}\def\@githubtitle{#2}}
\newcommand*{\githubsymbol}{\faGithub~}

% 2. Patch the title box to inject \github exactly the same way as \homepage
\patchcmd{\makecvtitle}
  {\ifthenelse{\isundefined{\@extrainfo}}}
  {%
    \ifthenelse{\isundefined{\@github}}{}{%
      \ifthenelse{\equal{\@githubtitle}{}}%
        {\makenewline\githubsymbol\link{\@github}}%
        {\makenewline\githubsymbol\link[\@githubtitle]{\@github}}}%
    \ifthenelse{\isundefined{\@extrainfo}}%
  }
  {}{}
\makeatother


\usepackage{lipsum} % Used for inserting dummy 'Lorem ipsum' text into the template
\usepackage[utf8]{inputenc}
\usepackage[margin=0.5in, bottom=0.8in]{geometry} % Increased bottom margin

\usepackage{fontspec}

\newfontface\lserif{Liberation Serif}
\newcommand{\Csh}{C{\lserif\#}}
%\setlength{\hintscolumnwidth}{3cm} % Uncomment to change the width of the dates column
\setlength{\makecvtitlenamewidth}{10cm} % For the 'classic' style, uncomment to adjust the width of the space allocated to your names
%----------------------------------------------------------------------------------------
%	NAME AND CONTACT INFORMATION SECTION
%---------------------------------------------s-------------------------------------------

\firstname{Zijian} % Your first name
\familyname{MENG} % Your last name

% All information in this block is optional, comment out any lines you don't need
%\title{Curriculum Vitae}
\address{269 Scout Street, Ottawa}{ON, Canada, K2C 4E7}
% \mobile{(+216) 52 573 153}
%\phone{(000) 111 1112}
%\fax{(000) 111 1113}
\email{contact@richardzjm.com}
% \email{richardzjm@gmail.com}
\homepage{richardzjm.com}{richardzjm.com} % The first argument is %the url for the clickable link, the second argument is the url displayed in the %template - this allows special characters to be displayed such as the tilde in this %example
\github{https://github.com/RichardZJM}{RichardZJM}
\photo[86pt][0pt]{ZJM_Logo.jpg} % The first bracket is the picture height, the second is %the thickness of the frame around the picture (0pt for no frame)
% \quote{ \Large Curriculum Vitae/Resume, Applying for PhD Mechanical Engineering}

%----------------------------------------------------------------------------------------

\begin{document}

\makecvtitle % Print the CV title

%----------------------------------------------------------------------------------------
%	EDUCATION SECTION
%----------------------------------------------------------------------------------------

\section{Education}
\cventry{Starting 2026}{Ph.D., Mechanical Engineering and Scientific Computing}{}{}{}{University of Michigan, Ann Arbor, USA}
\cventry{ Sep 2023 -- Dec 2025}{Master of Applied Science, Mechanical and Materials Engineering (Co-supervised with Electrical and Computer Engineering)}{}{}{}{Supervised by Laurent Karim Béland and Ryan Grant\\ Queen's University, Kingston, Canada\\4.3/4.3 GPA}  % Arguments not required can be left empty
\cventry{Sep 2018 -- May 2023}{Bachelor of Applied Science, Mechanical Engineering  }{}{}{}{With Professional Internship\\ Queen's University, Kingston, Canada\\4.18/4.3 GPA}


%----------------------------------------------------------------------------------------
%	WORK EXPERIENCE SECTION
%----------------------------------------------------------------------------------------

\section{Employment}

\cventry{Sep 2025 -- Dec 2025}{Computational Modelling Intern}{Canadian Nuclear Laboratories}{\small Chalk River,  Canada}{}{• Collaborated with experimental scientists to model helium in neutron-irradiated 304 stainless steel from a salvaged component of the National Research Universal reactor. \\ • Applied Object Kinetic Monte Carlo simulations for helium and defect distribution and effects.
}

\cventry{Fall 2023 \& 2024}{Lead Lab Teaching Assistant}{Queen's University}{\small Kingston, Canada}{}{• Engineering Graphics Course: Led 2 other TAs in lab sessions of more than 80 students. \\ •  Live demonstrations, management of sessions and students, tutorials, and marking.}

\cventry{Winter 2023}{Teaching Assistant}{Queen's University}{\small Kingston, Canada}{}{• Automatic Control Course: Communicating with students, marking, and proctoring.}

\cventry{May 2021 -- June 2022}{Product Engineering Intern}{Hanon Systems Canada}{\small Belleville, Canada}{}{• Developed and organized the builds of fluid transport prototypes for automotive thermals. \\
• Work included projects for electric vehicles from Ford, GM, and Lucid. \\
• Completed 342 prototype build requests under historically high demand. \\}



\section{Research Experience}
\cventry{Feb 2026 -- Present}{Research Assistant}{Queen's University}{\small Kingston, Canada}{}{Nuclear Materials Group, Computational Materials \\ • Training and running MLIPs for various projects and supporting software advances. \\ • Further improved MTP models, yielding up to 8$\,\times$  faster models with 3$\,\times$ lower error.} 

\cventry{Sep 2023 -- Dec 2025}{Master's Student Researcher}{Queen's University}{\small Kingston, Canada}{}{Nuclear Materials Group, Computational Materials \\ • Development, training, and validation of machine learning interatomic potentials in systems including Na–K–Cl–O, Ni, Si–O, Si–C; involved with Cu–H, Fe–Ni–Cr, Zr–H.  \\ • Proposed a cost-aware basis set pruning algorithm, optimizing cost-accuracy significantly. \\ •  Programmed a more optimized CPU and new GPU implementation of the Moment Tensor Potential in LAMMPS. \\ • Contributions yielded up to 13$\,\times$ CPU inference speedups at equal or better accuracy. } 
\cventry{Winter 2023}{Undergraduate Research Project}{Queen's University}{\small Kingston, Canada}{}{Nuclear Materials Group, Computational Materials \\ • Development, training, and validation of machine learning interatomic potentials.}

\section{Publications}
\cvitem{Published 2026}{\textbf{Meng, Z.}; Zongo, K.; Thoms, M.; Maxwell, C.; Torres, E.; Grant, R.; Béland, L. K. Accelerating moment tensor potentials through post-training pruning. Computational Materials Science 2026, 269, 114725. (Editor's Choice of the Month)}
\cvitem{Published 2026}{\textbf{Meng, Z.}; Zongo, K.; Torres, E.; Maxwell, C.; Grant, R.; Béland, L. K. A Kokkos-accelerated moment tensor potential implementation for LAMMPS. SoftwareX 2026, 33, 102524.}
\cvitem{Published 2025}{\textbf{Meng, Z.}; Sun, H.; Torres, E.; Maxwell, C.; Grant, R. E.; Béland, L. K. Small-cell-based fast active learning of machine learning interatomic potentials. Computational Materials Science 2025, 256, 113919.}


\section{Honors and Awards}
\cventry{2026}{Michigan Institute of Computational Discovery and Engineering Fellowship}{}{}{}{}
\cventry{2024}{UNENE Best Student Presentation}{}{}{}{Best student presentation at the UNENE 2024 conference in the 7-minute category.}
\cventry{2024}{NSERC Alexander Bell CGS-M Scholarship}{}{}{}{The Natural Sciences and Engineering Research Council of Canada's nationwide Master's graduate scholarship.}
\cventry{2023}{R. Samuel McLaughlin Fellowship}{}{}{}{}
\cventry{2023}{L. M. Arkley Prize}{}{}{}{Awarded for the best fourth-year Mechanical Engineering paper and presentation.}
\cventry{2023}{Colin T. Bayne Memorial Award}{}{}{}{Awarded to the graduating Mechanical Engineering student who, in the opinion of the Department, has shown most proficiency in innovative design.}
\cventry{2023}{Conn-Gilbert Award}{}{}{}{Awarded to the graduating Mechanical Engineering student with the highest average in core Thermodynamics courses.}
\cventry{2023}{3rd Place—Professional Engineers Ontario, Kingston, Engineering Competition}{}{}{}{}
\cventry{Various}{Dean's Scholar}{}{}{}{}
\cventry{2022}{Lorne C. Elder Scholarship}{}{}{}{}
\cventry{2021}{Best Project and Best Game Award, CuHacking Hackathon}{}{}{}{}
\cventry{2018}{Lena MacNeil Scholarship}{}{}{}{}

\section{Presentations and Posters}
\cventry{Presentation 2025}{Efficiency Advances in Atomistic Modelling of Nuclear Materials.}{}{}{}{UNENE 2025, Technical Presentation Category}
\cventry{Presentation 2025}{Moment Tensor Potential Efficiency Advances}{}{}{}{pyKMC 2025}
\cventry{Presentation 2024}{
Simulating Complex Chemistry for Advanced Reactors: Machine Learning Potentials in Action}{}{}{}{UNENE 2024, 7-minute Category}
\cventry{Poster \\2023}{Systematic Small-Cell Training Set Selection}{}{}{}{UNENE Student Posters}







%----------------------------------------------------------------------------------------
%	LANGUAGES SECTION
%----------------------------------------------------------------------------------------

% \section{Langues}
% \begin{small}
% \cvitemwithcomment{Arabe}{Maternelle\\dhe}{}
% \cvitemwithcomment{Français}{Très Bon Niveau}{DELF B1 - 2008}
% \cvitemwithcomment{Anglais}{Très Bon Niveau}{TOEIC B2 - 2016}
% \end{small}


% %----------------------------------------------------------------------------------------
% %	INTERESTS SECTION
% %----------------------------------------------------------------------------------------
% %\bigskip

% \section{Centres d'intérêt }

% \renewcommand{\listitemsymbol}{-~} % Changes the symbol used for lists

% \cvitem{Club} {Vice Président Club Microsoft ENETCOM (2014)} 
% \cvitem{Loisirs} {Voyage, Sport, Musique } 


%----------------------------------------------------------------------------------------
%	COVER LETTER
%----------------------------------------------------------------------------------------

% To remove the cover letter, comment out this entire block

%\clearpage

%\recipient{HR Departmnet}{Corporation\\123 Pleasant Lane\\12345 City, State} % Letter recipient
%\date{\today} % Letter date
%\opening{Dear Sir or Madam,} % Opening greeting
%\closing{Sincerely yours,} % Closing phrase
%\enclosure[Attached]{curriculum vit\ae{}} % List of enclosed documents

%\makelettertitle % Print letter title

%\lipsum[1-3] % Dummy text

%\makeletterclosing % Print letter signature

%----------------------------------------------------------------------------------------

\end{document}